\begin{textsample}{POS Dim 2 – 2019\_c4 – Score 18.00 – 2019\_c4/enade\_2019\_biomed\_ef38a0b72c.txt}  \label{ex:f2_pos_041}
CAPÍTULO 4 PERCEPÇÃO DA PROVA Nas análises feitas neste capítulo, constam as percepções dos concluintes da Área de Biomedicina sobre a prova aplicada no Enade/2019.
Essas percepções foram mensuradas por meio de nove questões que avaliaram desde o grau de dificuldade da prova até o tempo gasto para concluí -la.
As percepções sobre a prova foram relacionadas ao desempenho dos estudantes, à Grande Região de funcionamento do curso, à Categoria Administrativa e ao tipo de Organização Acadêmica da IES.
O \textbf{questionário} de percepção da prova encontra -se ao final do Anexo VII, onde está a \textbf{reprodução} do exame.
O desempenho dos estudantes foi classificado em quatro quartos.
Para tanto, esse desempenho foi \textbf{ordenado} de forma ascendente.
O percentil 25, P25, também conhecido como primeiro quartil, é a nota de desempenho que deixa um quarto ( 25 \% ) dos valores observados abaixo e três quartos acima.
Na Figura 1, apresenta -se uma ilustração desse conceito.
O quarto inferior de desempenho é composto pelas notas abaixo do primeiro quartil.
Já o percentil 75, P75, também conhecido como terceiro quartil, é o valor para o qual há três quartos ( 75 \% ) dos dados abaixo e um quarto acima dele.
O quarto superior de desempenho é composto pelas notas iguais ou acima do terceiro quartil.
O percentil 50, P50, também conhecido como mediana, é o valor que divide as notas em dois conjuntos de igual tamanho.
O segundo quarto inclui valores entre o primeiro quartil ( P25 ) e a mediana.
O terceiro quarto contém os valores entre a mediana ( P50 ) e o terceiro quartil ( P75 ).
Vale ressaltar que percentis, quartis e medianas são pontos que, não obrigatoriamente, pertencem ao conjunto original de dados, ao passo que os quartos são \textbf{subconjuntos} dos dados originais.
A seguir, serão apresentados gráficos com resultados selecionados, relativos às nove questões avaliadas por grupos de estudantes.
Nas barras dos gráficos, apresenta -se o percentual de alunos que assinalou uma das opções ou a soma das porcentagens daqueles que assinalaram duas ( ou três ) delas.
Por exemplo, para as questões 1 e 2, nos gráficos, é apresentada a porcentagem total de participantes que assinalou as opções ( D ) Difícil e ( E ) Muito difícil.
Em cada barra, foram assinalados, também, os extremos do intervalo de confiança de 95 \% como linhas verticais unidas por uma linha horizontal na forma da letra H maiúscula.
O \textbf{estimador} de um \textbf{parâmetro} com um certo nível de confiança ( e.g. 95,0 \% ) deve conter o \textbf{parâmetro} no intervalo de confiança em 95 \% das vezes.
Na comparação entre os \textbf{estimadores} dos \textbf{parâmetros} de duas classes de uma dada categoria ( e.g.
Norte e Nordeste nas Grandes Regiões, de primeiro e último quarto dentro de desempenho, de IES Públicas e Privada, ou de Universidades e Faculdades ), associados aos seus respectivos intervalos de confiança, diz -se que não há diferença estatisticamente significativa entre os \textbf{parâmetros} das duas categorias quando há \textbf{interseção} entre os intervalos de confiança e que há diferença, se os intervalos de confiança são \textbf{disjuntos} ( para maiores informações \textbf{vide} Glossário ).
Nas Tabelas do Anexo II, são apresentados os valores absolutos e a distribuição percentual18 das alternativas válidas das nove questões, segundo o mesmo \textbf{recorte} de desempenho dos alunos e da Grande Região de funcionamento do curso, de Categoria Administrativa e do tipo de Organização Acadêmica da IES. 18 Cumpre lembrar uma das convenções para tabelas \textbf{numéricas} ( pág. iii ) sobre a possibilidade de a soma das partes não resultar em 100 \% por questões de \textbf{arredondamento}. 4.1 GRAU DE DIFICULDADE DA PROVA 4.1.1 Componente de Formação Geral Ao avaliarem Qual o grau de dificuldade desta prova na parte de Formação Geral? ( Questão 1 ), 30,6 \% do grupo de inscritos e presentes optaram pelas alternativas Difícil ou Muito difícil.
Entretanto, para mais da metade dos estudantes ( 62,0 \% ), o Componente de Formação Geral da prova foi avaliado com grau de dificuldade Médio ( Gráfico 4.1, Gráfico 4.2 e, no Anexo II, a Tabela II.1 ).
O percentual de estudantes que considerou a prova Difícil ou Muito difícil foi maior na região Centro-Oeste, onde a proporção foi de 36,0 \%, enquanto a de menor proporção foi a Nordeste, com 27,6 \%.
No Gráfico 4.1, é possível observar que as diferenças entre a região Nordeste e as regiões Sudeste e Sul são estatisticamente significativas, assim como a diferença entre a região Centro-Oeste e as regiões Norte, Nordeste e Sudeste.
Nas Grandes Regiões, a proporção de presentes à prova que avaliou o Componente de Formação Geral como de grau de dificuldade Médio esteve entre 58,4 \%, na região Centro-Oeste, e 64,2 \%, na região Norte.
O percentual de alunos que considerou o componente de Formação Geral da prova Difícil ou Muito difícil foi decrescente em função dos quartos de desempenho : 34,6 \%, no primeiro quarto, e 28,3 \%, no quarto, o grupo de melhor desempenho na prova.
Nos quartos de desempenho intermediários, a proporção de alunos que considerou a prova Difícil ou Muito difícil foi de 30,4 \%, no segundo quarto, e de 29,0 \%, no terceiro quarto.
As diferenças entre o quarto inferior de desempenho e os demais são estatisticamente significativas.
Para todos os quartos de desempenho, a alternativa modal para essa pergunta foi a Médio, com 59,7 \% e 62,6 \% dos respondentes nos quartos extremos, primeiro e quarto, respectivamente. como de grau de dificuldade Médio esteve entre 58,4 \%, na região Centro-Oeste, e 64,2 \%, na região Norte.
O percentual de alunos que considerou o componente de Formação Geral da prova Difícil ou Muito difícil foi decrescente em função dos quartos de desempenho : 34,6 \%, no primeiro quarto, e 28,3 \%, no quarto, o grupo de melhor desempenho na prova.
Nos quartos de desempenho intermediários, a proporção de alunos que considerou a prova Difícil ou Muito difícil foi de 30,4 \%, no segundo quarto, e de 29,0 \%, no terceiro quarto.
As diferenças entre o quarto inferior de desempenho e os demais são estatisticamente significativas.
Para todos os quartos de desempenho, a alternativa modal para essa pergunta foi a Médio, com 59,7 \% e 62,6 \% dos respondentes nos quartos extremos, primeiro e quarto, respectivamente.
O percentual de estudantes que considerou o componente de Formação Geral da prova Difícil ou Muito difícil foi maior entre os estudantes de Faculdades, alcançando 32,9 \%, enquanto entre os estudantes das Universidades, a Organização Acadêmica com a menor incidência, alcançou 28,9 \%.
No Gráfico 4.3, é possível observar que as diferenças entre Faculdades e Universidades são estatisticamente significativas.
Nas Organizações Acadêmicas, a proporção de presentes à prova que avaliou o Componente de Formação Geral como de grau de dificuldade Médio esteve entre 61,6 \%, nas Faculdades e Centros Universitários, e 62,6 \% nas Universidades ( ver também a Tabela II.2, no Anexo II ).
O percentual de estudantes que considerou o componente de Formação Geral da prova Difícil ou Muito difícil foi maior entre os estudantes de cursos de IES Privadas ( 31,2 \% ), com diferença estatisticamente significativa da proporção de IES Públicas ( 24,5 \% ).
Já a proporção de presentes à prova que avaliou o Componente de Formação Geral como de grau de dificuldade Médio foi de 61,9 \% para os de IES Públicas e de 62,0 \% para os de IES Privadas ( ver Gráfico 4.4 ). 4.1.2 Componente de Conhecimento Específico Ao responderem à Questão 2 – Qual o grau de dificuldade desta prova na parte de Componente Específico? – 37,6 \% do grupo de estudantes classificaram -na como Difícil ou Muito difícil.
Além disso, o Componente modal de Conhecimento Específico da prova foi considerado de grau de dificuldade Médio por 58,3 \% dos alunos ( Gráfico 4.5, Gráfico 4.6, e, no Anexo II, a Tabela II.3 ).
A análise das respostas dos estudantes quanto ao grau de dificuldade do Componente de Conhecimento Específico da prova, agregado por Grande Região, mostra diferenças estatisticamente significativas entre proporções de alunos que a avaliaram como Difícil ou Muito difícil.
Tais diferenças são observadas entre a região com a maior proporção, a Sudeste ( 41,0 \% ), e as regiões Norte ( 35,2 \% ), Nordeste ( 32,0 \% ) e Centro-Oeste ( 33,5 \% ).
As diferenças entre as regiões Nordeste ( 32,0 \% ) e Sul ( 37,9 \% ) também são estatisticamente significativas.
O percentual de alunos que classificou o grau de dificuldade como Médio, no Componente de Conhecimento Específico, variou de 55,6 \% a 63,0 \%, nas regiões Sudeste e Nordeste, respectivamente.
Considerando -se a avaliação da dificuldade das questões do Componente de Conhecimento Específico da prova de acordo com o desempenho dos estudantes, observa -se uma tendência decrescente dos que consideraram a parte específica Difícil ou Muito difícil, variando de 30,8 \% ( último quarto ) a 42,8 \% ( primeiro quarto ).
As demais proporções foram de 39,9 \% e 37,0 \%, respectivamente, no segundo e terceiro quartos.
Observa -se que há diferença estatisticamente significativa dos resultados entre o quarto superior de desempenho e os demais, assim como entre o quarto inferior e o terceiro quarto.
Já os que responderam que o grau de dificuldade das questões do Componente de Conhecimento Específico da prova foi Médio atingiram \textbf{percentagens} que variaram de 53,3 \%, no primeiro quarto, a 63,8 \%, no último.
Considerando -se a avaliação da dificuldade das questões do Componente de Conhecimento Específico da prova, o percentual de estudantes que as considerou Difícil ou Muito difícil foi maior entre os estudantes de Faculdade, atingindo 41,9 \% ; já entre os de Universidades, a Organização Acadêmica com a menor incidência, essa proporção foi de 36,1 \%.
No Gráfico 4.7, é possível observar que a diferença entre Faculdades e Universidades é estatisticamente significativa, bem como a entre Faculdades e Centros Universitários ( 36,7 \% ).
Nas Organizações Acadêmicas, a proporção de presentes à prova que considerou o Componente de Conhecimento Específico de grau de dificuldade Médio esteve entre 53,9 \%, nas Faculdades, e 60,2 \%, nas Universidades ( ver também Tabela II.4, no Anexo II ).
O percentual de estudantes que considerou o Componente Específico da prova Difícil ou Muito difícil foi maior entre os estudantes de cursos de IES Privadas ( 38,0 \% ), com diferença estatisticamente significativa da proporção em IES Públicas ( 33,5 \% ).
Já a proporção de presentes à prova que avaliou este componente como de grau de dificuldade Médio foi de 62,0 \% para os de IES Públicas e de 57,9 \% para os de IES Privadas ( ver Gráfico 4.8 e Tabela II.4, no Anexo II ). 4.2 EXTENSÃO DA PROVA EM RELAÇÃO AO TEMPO TOTAL \textbf{Indagados} quanto à extensão da prova, em relação ao tempo total oferecido para a sua resolução ( Questão 3 ), os estudantes apontaram, com maior incidência, a alternativa que considerava a extensão adequada, para todas as agregações consideradas ( Gráfico 4.9, Gráfico 4.10, e, no Anexo II, a Tabela II.5 ).
O percentual de alunos que respondeu ser a extensão da prova adequada foi de 67,8 \%.
Já 26,3 \% dos inscritos presentes consideraram que a prova foi Longa ou Muito longa, e 6,0 \% a avaliaram como Curta ou Muito curta.
Dentre as Grandes Regiões, os que avaliaram a prova como Longa ou Muito longa, em relação ao tempo total destinado à sua resolução, atingiram uma porcentagem que variou de 20,3 \%, na região Norte, até 27,8 \%, na região Sudeste.
Há diferença estatisticamente significativa entre as regiões Sudeste e Sul, e também entre a região Norte e as regiões Nordeste e Sudeste.
Considerando -se o desempenho dos alunos, nota -se ainda que 64,1 \% consideraram a extensão da prova adequada, no quarto de desempenho inferior, e 69,3 \%, no de melhor desempenho ( quarto superior ).
Nos quartos intermediários, essa proporção foi de 68,1 \%, no segundo quarto, e de 69,4 \%, no terceiro.
No Gráfico 4.10, observa -se que as proporções de estudantes que consideraram a prova Longa ou Muito longa apresentam tendência decrescente em função dos quartos de desempenho, sendo a maior proporção dessa resposta no primeiro quarto de desempenho ( 29,9 \% ), e a menor, no quarto superior ( 24,1 \% ).
Pode -se constatar que há diferença estatisticamente significativa entre a proporção do quarto inferior de desempenho e a dos demais.
Considerando -se a extensão da prova, o percentual de estudantes que avaliou a prova como Longa ou Muito longa, em relação ao tempo total destinado à sua resolução, foi maior entre os estudantes de Centros Universitários, alcançando 27,3 \%, enquanto entre os estudantes de Universidades, a Organização Acadêmica com a menor incidência, esse índice atingiu 25,4 \%, e entre os de Faculdades, 26,3 \%.
No Gráfico 4.11, é possível observar que não há diferença estatisticamente significativa entre os três tipos de Organização Acadêmica.
A proporção de presentes à prova que considerou a sua extensão Adequada, em relação ao tempo total destinado à sua resolução, esteve entre 66,4 \%, nas Faculdades, e 69,2 \%, nas Universidades. ( ver também Tabela II.6, no Anexo II ).
O percentual de estudantes que considerou a prova Longa ou Muito longa, em relação ao tempo total destinado à sua resolução, foi maior entre os estudantes de cursos de IES Privadas ( 26,6 \% ), não havendo diferença estatisticamente significativa da proporção em relação aos de IES Públicas ( 23,0 \% ).
Já a proporção de presentes à prova que considerou a extensão Adequada foi de 71,7 \% para os de IES Públicas e de 67,4 \% para os de IES Privadas ( ver Gráfico 4.12 e Tabela II.6, no Anexo II ). 4.3 COMPREENSÃO DOS ENUNCIADOS DAS QUESTÕES 4.3.1 Componente de Formação Geral Em relação aos enunciados das questões do Componente de Formação Geral ( Questão 4 ), as opiniões foram positivas, já que 75,2 \% dos alunos avaliados consideraram todos ou a maioria dos enunciados das questões claros e objetivos ( Gráfico 4.13, Gráfico 4.14 e, no Anexo II, a Tabela II.7 ).
Na análise regional, a \textbf{percentagem} de estudantes que avaliou que todos ou a maioria dos enunciados das questões do Componente de Formação Geral estavam claros e objetivos variou de 71,9 \%, na região Nordeste, a 78,9 \%, na região Sul.
A análise das percepções dos estudantes sobre a clareza e a objetividade dos enunciados permite afirmar que todos, ou a maioria dos enunciados de questões relativas ao Componente de Formação Geral, foram considerados claros e objetivos para a maior parte dos respondentes ( maior ou igual a 71,9 \%, em todas as regiões e maior ou igual a 71,8 \% para todos os quartos de desempenho ).
Segundo o desempenho, observa -se que a proporção dos que emitiram essa opinião \textbf{cresce}, conforme o desempenho aumenta, com diferenças estatisticamente significativas entre o quarto inferior de desempenho e os demais.
No quarto superior, a clareza e a objetividade de todos ou da maioria dos enunciados das questões foram percebidas por 77,2 \%.
A proporção dos presentes que considerou que todos ou a maioria dos enunciados das questões do Componente de Formação Geral estavam claros e objetivos foi maior entre os estudantes de Centros Universitários, chegando a 75,7 \%, enquanto entre os de Faculdades, a Organização Acadêmica com a menor incidência, esse índice atingiu 74,0 \%, e entre os de Universidades, 75,4 \%.
No Gráfico 4.15, é possível observar que as diferenças entre os três tipos de Organização Acadêmica não são estatisticamente significativas ( ver também Tabela II.8, no Anexo II ).
O percentual de estudantes que considerou que todos ou a maioria dos enunciados das questões do Componente de Formação Geral estavam claros e objetivos foi maior entre os estudantes de cursos de IES Públicas ( 75,3 \% ), mas sem diferença estatisticamente significativa da proporção em IES Privadas ( 75,2 \% ).
Um maior detalhamento pode ser encontrado no Gráfico 4.16 e na Tabela II.8, no Anexo II. 4.3.2 Componente de Conhecimento Específico Em relação aos enunciados das questões do Componente de Conhecimento Específico da prova, para 85,7 \% dos estudantes avaliados da Área de Biomedicina, a clareza e a objetividade ( Questão 5 ) estavam presentes em todas ou na maioria das questões ( Gráfico 4.17 a Gráfico 4.20, e no Anexo II, as Tabelas II.9 e II.10 ).
A maioria dos estudantes de todas as Grandes Regiões brasileiras considerou claros e objetivos todos ou a maioria dos enunciados das questões do Componente de Conhecimento Específico da prova, percentual sempre maior ou igual a 82,3 \%.
A diferença entre a região Nordeste e as regiões Sudeste e Centro-Oeste é estatisticamente significativa, bem como entre as regiões Norte e Centro-Oeste e entre a região Sul e as regiões Norte, Nordeste e Sudeste.
A proporção de estudantes que considerou os enunciados das questões claros e objetivos apresenta uma tendência crescente em relação ao aumento de desempenho : mais elevada no quarto superior ( 91,9 \% ), se comparada ao quarto inferior de desempenho ( 76,8 \% ).
As diferenças entre todos os quartos de desempenho são estatisticamente significativas.
A proporção dos presentes que considerou que todos ou a maioria dos enunciados das questões do Componente de Conhecimento Específico estavam claros e objetivos foi maior entre os estudantes de Universidades, atingindo 87,0 \%.
Entre os estudantes de Faculdades, a Organização Acadêmica com a menor incidência, essa proporção atingiu 81,8 \%, sendo intermediária a proporção de Centro Universitários ( 86,4 \% ).
No Gráfico 4.19, é possível observar que as diferenças entre Faculdades e os demais tipos de Organização Acadêmica são estatisticamente significativas ( ver também Tabela II.10, no Anexo II ).
O percentual de estudantes que considerou que todos ou a maioria dos enunciados das questões do Componente de Conhecimento Específico estavam claros e objetivos foi maior entre os estudantes de cursos de IES Públicas ( 89,6 \% ), com diferença estatisticamente significativa da proporção em IES Privadas ( 85,3 \% ).
Ver Gráfico 4.20 e Tabela II.10, no Anexo II, para maior detalhamento. 4.4 SUFICIÊNCIA DAS INFORMAÇÕES/INSTRUÇÕES FORNECIDAS Ao se avaliarem as informações/instruções fornecidas para a resolução das questões ( Questão 6 ), constatou -se que 84,7 \% dos respondentes da Área de Biomedicina de todo o Brasil afirmaram que estas eram até excessivas ou suficientes em todas ou na maioria das questões ( Gráfico 4.21, Gráfico 4.22 e, no Anexo II, a Tabela II.11 ).
Quanto à distribuição de respondentes pelas Grandes Regiões, observa -se que a proporção de estudantes que considerou as informações/instruções fornecidas até excessivas ou suficientes em todas ou na maioria das questões foi sempre superior ou igual a 81,8 \%, chegando a 88,9 \%, na região Sul.
Existem diferenças estatisticamente significativas entre as regiões Sudeste e Sul e entre a região Nordeste e as regiões Sul e Centro-Oeste.
Levando -se em conta o desempenho dos participantes, observa -se uma tendência crescente das proporções de participantes que avaliaram as informações/instruções como até excessivas ou suficientes em todas ou na maioria das questões, havendo diferenças estatisticamente significativas entre o primeiro, o segundo e o quarto superior, bem como entre o terceiro e o quarto superior de desempenho, conforme mostrado no Gráfico 4.22.
O percentual foi mais elevado no quarto superior ( 88,8 \% ), percentual superior à média nacional ( 84,7 \% ).
Já no quarto inferior, a suficiência das informações/instruções declarada como até excessiva, em todas ou na maioria das questões foi percebida por 80,3 \% dos respondentes.
Considerando -se as informações/instruções fornecidas para a resolução das questões ( Questão 6 ), a proporção que considerou que estas eram até excessivas ou suficientes em todas ou na maioria das questões foi maior entre os estudantes de Universidades, \textbf{totalizando} 86,1 \%.
As Faculdades foram a Organização Acadêmica em que houve a menor proporção, 81,4 \%.
No Gráfico 4.23, é possível observar que as diferenças entre Faculdades ( 81,4 \% ), em relação aos Centros Universitários ( 85,0 \% ) e às Universidades ( 86,1 \% ), são estatisticamente significativas ( ver também Tabela II.12, no Anexo II ).
O percentual de estudantes que considerou que as informações/instruções fornecidas para a resolução das questões eram até excessivas ou suficientes em todas ou na maioria das questões foi maior para os estudantes de cursos de IES Públicas ( 87,1 \% ), não havendo diferença estatisticamente significativa da proporção em IES Privadas ( 84,5 \% ). ( ver Gráfico 4.24 e Tabela II.12, no Anexo II ). 4.5 DIFICULDADE ENCONTRADA AO RESPONDER À PROVA \textbf{Indagados} sobre as dificuldades com as quais se depararam ao responder à prova ( Questão 7 ), 20,6 \% dos estudantes apontaram o Desconhecimento do conteúdo.
Para 54,6 \%, a Forma diferente de abordagem do conteúdo foi indicada como dificuldade.
Já a Falta de motivação para fazer a prova foi a dificuldade apontada por 7,0 \% dos respondentes.
Considerando -se todo o Brasil, 15,3 \% dos respondentes afirmaram que não tiveram qualquer tipo de dificuldade para responder à prova ( Tabelas II.13 e II.14, no Anexo II ).
Nos Gráficos 4.25 a 4.28, são apresentados os percentuais de estudantes que apontaram o Desconhecimento do conteúdo como dificuldade ao responder à prova.
Na análise por Grandes Regiões, o percentual de inscritos e presentes que apontaram o Desconhecimento do conteúdo como dificuldade ao responder à prova não superou 21,5 \%.
Os percentuais variaram de 19,0 \%, na região Nordeste, a 21,5 \%, na região Sul.
A Forma diferente de abordagem do conteúdo foi a escolha modal dos estudantes, com percentuais que variaram de 50,9 \% ( região Centro-Oeste ) a 57,1 \% ( região Norte ).
O percentual de alunos que citou a Falta de motivação como dificuldade variou de 5,8 \% ( região Norte ) a 8,4 \% ( região Nordeste ).
Os que declararam não ter qualquer tipo de dificuldade para responder à prova variaram de 13,2 \%, na região Norte, a 19,5 \%, na região Centro-Oeste.
Em relação aos quartos de desempenho, o Desconhecimento do conteúdo foi a opção escolhida por 18,2 \% dos estudantes do quarto inferior e por 22,3 \% do quarto superior, caracterizando uma tendência crescente.
São observadas diferenças estatisticamente significativas entre o quarto inferior e os dois quartos superiores de desempenho.
A alternativa modal para os alunos, quando agregados pelos quartos de desempenho, para a dificuldade encontrada, foi causada pela Forma diferente de abordagem do conteúdo : 57,2 \%, no primeiro quarto, e 57,4 \% do segundo quarto, escolheram essa alternativa.
Na análise por Organização Acadêmica, o percentual de inscritos e presentes que apontou o Desconhecimento do conteúdo como dificuldade ao responder à prova não superou 20,9 \%.
Os percentuais variaram de 20,1 \%, nas Faculdades, a 20,9 \%, nas Universidades.
A Forma diferente de abordagem do conteúdo foi a escolha modal dos estudantes, com percentuais que variaram de 52,9 \% ( Universidades ) a 58,0 \% ( Faculdades ).
O percentual de alunos que citou a Falta de motivação como dificuldade variou de 5,8 \% ( Centros Universitários ) a 8,3 \% ( Universidades ).
Os que declararam não ter qualquer tipo de dificuldade para responder à prova variaram de 12,9 \%, nas Faculdades, a 16,4 \%, nos Centros Universitários.
No Gráfico 4.27, é possível observar que não há diferenças estatisticamente significativas entre Faculdades, Centros Universitários e Universidades ( ver também Tabela II.14, no Anexo II ).
Em relação às Categorias Administrativas, o Desconhecimento do conteúdo foi a opção escolhida por 27,7 \% dos estudantes nas IES Públicas e por 19,9 \% nas IES Privadas.
São observadas diferenças estatisticamente significativas entre essas duas categorias.
A alternativa modal para os alunos, quando agregados pelas Categorias Administrativas, para a dificuldade encontrada, foi causada pela Forma diferente de abordagem do conteúdo : 42,8 \%, para IES Públicas e 55,8 \%, para as IES Privadas. ( ver Gráfico 4.28 e Tabela II.14, no Anexo II ). 4.6 CONTEÚDOS DAS QUESTÕES OBJETIVAS DA PROVA Ao se analisarem os conteúdos das questões objetivas da prova ( Questão 8 ), verifica -se que um percentual muito pequeno dos estudantes avaliados, apenas 2,4 \%, afirmou que não estudou ainda a maioria desses conteúdos ( Gráficos 4.29 a Gráfico 4.32, e as Tabelas II.15 e II.16, no Anexo II ).
A maioria ( 72,8 \% ) afirmou ter estudado e aprendido muitos ou todos os conteúdos avaliados.
Na análise por Grande Região, a proporção de respondentes que escolheu a opção não estudou ainda a maioria desses conteúdos foi pequena.
Observa -se que, nas regiões Norte ( 3,1 \% ) e Nordeste ( 2,9 \% ), as proporções foram ligeiramente maiores do que a média nacional ( 2,4 \% ).
Nas regiões Sul e Centro-Oeste, a situação foi \textbf{inversa}, com proporções mais baixas, correspondendo a, respectivamente, 1,2 \% e 1,9 \%.
Na região Sudeste, a proporção de respondentes que escolheram essa resposta foi a mesma da nacional ( 2,4 \% ).
São observadas diferenças estatisticamente significativas entre a região Sul e as regiões Norte, Nordeste e Sudeste.
Em todas as regiões, a maioria dos presentes afirmou ter estudado e aprendido muitos ou todos os conteúdos, com proporções variando entre 68,4 \%, na região Norte, a 81,2 \%, na região Sul.
Considerando -se separadamente as opiniões de estudantes dos quatro quartos de desempenho, observa -se que, no quarto inferior, 4,6 \% ofereceram como resposta que não estudou ainda a maioria desses conteúdos, sendo 0,9 \% os do quarto superior com a mesma resposta.
A diferença entre os alunos que optaram por esse motivo de dificuldade no quarto inferior em relação aos demais quartos de desempenho é estatisticamente significativa, bem como a diferença entre os do segundo quarto e os do quarto de desempenho superior.
Tendo -se em conta o quarto superior, 85,2 \% dos alunos afirmaram ter estudado e aprendido muitos ou todos os conteúdos.
No outro extremo, no primeiro quarto, 58,8 \% optaram pelas mesmas categorias.
As proporções são crescentes com o desempenho.
Na análise por Organização Acadêmica, a proporção de respondentes que escolheu a opção não estudou ainda a maioria desses conteúdos foi pequena.
Observa -se que as proporções nos Centros Universitários e Universidades, ambos com 2,0 \%, foram menores do que a média nacional ( 2,4 \% ).
Nas Faculdades, a situação foi \textbf{inversa} : a proporção foi mais alta ( 3,8 \% ).
No Gráfico 4.31, é possível observar que as diferenças entre Faculdades ( 3,8 \% ) e os demais tipos de Organização Acadêmica são estatisticamente significativas.
Em todos os tipos de Organização Acadêmica, a maioria dos presentes afirmou ter estudado e aprendido muitos ou todos os conteúdos, com proporções variando entre 66,7 \%, nas Faculdades, a 75,5 \%, nas Universidades. ( ver também Tabela II.16, no Anexo II ).
Em relação às Categorias Administrativas, observa -se que, nas IES Públicas, 2,8 \% ofereceram como resposta que não estudou ainda a maioria desses conteúdos, sendo 2,4 \% nas IES Privadas com a mesma resposta.
A diferença entre os alunos que optaram por esse motivo de dificuldade nas duas categorias não é estatisticamente significativa.
Considerando -se as IES Públicas, 76,0 \% dos alunos afirmaram ter estudado e aprendido muitos ou todos os conteúdos, sendo de 72,5 \% a proporção dos que escolheram essa resposta nas IES Privadas. ( ver Gráfico 4.32 e Tabela II.16, no Anexo II ). 4.7 TEMPO GASTO PARA CONCLUIR A PROVA Ao responderem sobre o tempo de conclusão da prova ( Questão 9 ), mais de quatro quintos dos estudantes ( 86,4 \% ) afirmaram ter gasto Entre duas e quatro horas ( Gráfico 4.33 a Gráfico 4.36 e, no Anexo II, as Tabelas II.17 e II.18 ).
Considerando -se as cinco Grandes Regiões brasileiras, os percentuais dos que utilizaram Entre duas e quatro horas para finalizar a prova nas regiões Norte ( 84,7 \% ) e Nordeste ( 85,0 \% ) foram inferiores ao percentual nacional.
Nas demais Grandes Regiões, os percentuais de alunos que \textbf{dispensaram} Entre duas e quatro horas para concluir a prova ficaram muito próximos ; 87,2 \% na região Sudeste ; 86,9 \%, na região Sul ; e 86,6 \%, na região Centro-Oeste, como mostrado no Gráfico 4.33.
Não há diferenças estatisticamente significativas entre as regiões.
Uma vez analisadas as alternativas escolhidas pelos estudantes que se situam nos diferentes quartos de desempenho, observa -se uma tendência crescente : uma proporção maior de participantes no quarto superior declarou ter gasto Entre duas e quatro horas para concluir a prova quando comparada com a proporção de participantes nos quartos inferiores.
Há diferenças estatisticamente significativas entre o quarto inferior e os demais quartos de desempenho, bem como entre o segundo e o quarto superior de desempenho.
Na análise por Organização Acadêmica, os percentuais dos que utilizaram Entre duas e quatro horas para finalizar a prova nas Universidades ( 86,5 \% ) e Centros Universitários ( 86,9 \% ) foram superiores ao percentual nacional ( 86,4 \% ).
Nas Faculdades, tal percentual foi de 85,2 \%, e, como mostrado no Gráfico 4.35, não há diferença estatisticamente significativa entre as proporções dos três tipos de Organização Acadêmica.
Em relação às Categorias Administrativas, observa -se que, nas IES Públicas, 82,4 \% declararam ter gasto Entre duas e quatro horas para concluir a prova, sendo 86,7 \% nas IES Privadas com a mesma resposta.
A diferença entre os alunos que optaram por esse motivo de dificuldade nas duas categorias é estatisticamente significativa ( ver Gráfico 4.36 e Tabela II.18, no Anexo II ).

% matched lemmas: arredondamento, crescer, disjunto, dispensar, estimador, indagar, interseção, inverso, numérico, ordenar, parâmetro, percentagem, questionário, recorte, reprodução, subconjunto, totalizar, vidar
\end{textsample}
